\documentclass{article}
\usepackage[utf8]{inputenc}
\usepackage{graphicx}
\usepackage{amsmath}
\usepackage{amssymb}
\usepackage{url}
\usepackage{tabularx}
\usepackage[colorlinks=true,allcolors=blue]{hyperref}
\usepackage{caption}
\usepackage{pdfpages}
\usepackage{multicol}
\usepackage[ngerman]{babel}
\usepackage{multicol}
\usepackage{geometry}	% [textwidth=17cm,textheight = 24.5cm]
% Um die Dokumentränder zu sehen: [showframe] vor {geometry} einfügen
\usepackage{bbold}
\usepackage{makecell}
\usepackage{float}
\usepackage{braket}
\usepackage{cancel}
\usepackage{braket}
\usepackage{mathtools}
\usepackage{xltabular}
\usepackage{float}
\usepackage{MnSymbol}
\usepackage{multirow}
\usepackage{chemformula}

\renewcommand{\cellalign}{tl}

\setlength{\columnsep}{1cm}
\setlength\parindent{0pt}

\newcommand{\diff}{\mathrm{d}}
\newcommand{\konst}{\mathrm{konst.}}
\newcommand{\lr}[1]{\left(#1\right)}
\newcommand{\lreck}[1]{\left[#1\right]}
\newcommand{\dkon}[3]{\left(\frac{\partial #1}{\partial #2}\right)_{#3}}

\definecolor{gruen}{HTML}{038133}

\newcommand{\commenttext}[1]{\textcolor{gruen}{\text{#1}}}
\newcommand{\comment}[1]{{\color{gruen} {#1}}}

\begin{document}

\begin{center}
	\Large{\textbf{Theoretische Physik I: Analytische Mechanik}}\\
	\small{613348}\\
	\small{\today}
\end{center}

\tableofcontents


\newpage


\section*{Übersicht}
Die analytische Mechanik führt neben dem Newton'schen Formulismus neue Formalismen, die zur Beschreibung klassischer mechanischer Systeme dienen, sowie das Prinzip der kleinsten Wirkung ein.\\
Diese Konzepte sind in anderen Bereichen der Physik ebenso von großer Bedeutung.\\
\\
Mit dem Noether-Theorem gelingt die Verbindung zwischen den Symmetrien des betrachteten Systems und den Erhaltungsgrößen.\\
\\
Zuletzt werden das Zweikörperproblem sowie die Mechanik des starren Körpers mithilfe des Lagrange-Formalismus beschrieben.


\newpage


\section{Mathematische Grundlagen}
\subsection{Transformationen}
Die Rotationen um die $x-,y-,z-$Achse um den Winkel $\varphi$, sowie dessen Verbindungen mit den Erzeugenden $T_{x,y,z}$ lauten
	\begin{alignat}{2}
		\underline{A}_{x,\varphi} &=
		\begin{pmatrix}
			1 & 0 & 0 \\
			0 & \cos\varphi & \sin\varphi \\
			0 & -\sin\varphi & \cos\varphi
		\end{pmatrix}
		&& = \exp\lr{\varphi\underline{T}_x} = \exp\lreck{\varphi 
		\begin{pmatrix}
			0 & 0 & 0 \\
			0 & 0 & 1 \\
			0 & -1 & 0
		\end{pmatrix}} \\
		\underline{A}_{y,\varphi} &=
		\begin{pmatrix}
			\cos\varphi  & 0 & -\sin\varphi \\
			0 & 1 & 0 \\
			\sin\varphi & 0 & \cos\varphi
		\end{pmatrix}
		&& = \exp\lr{\varphi\underline{T}_y} = \exp\lreck{\varphi 
		\begin{pmatrix}
			0 & 0 & -1 \\
			0 & 0 & 0 \\
			1 & 0 & 0
		\end{pmatrix}} \\
		\underline{A}_{z,\varphi} &=
		\begin{pmatrix}
			\cos\varphi & \sin\varphi & 0 \\
			-\sin\varphi & \cos\varphi & 0 \\
			0 & 0 & 1
		\end{pmatrix}
		&& = \exp\lr{\varphi\underline{T}_z} = \exp\lreck{\varphi 
		\begin{pmatrix}
			0 & 1 & 0 \\
			-1 & 0 & 0 \\
			0 & 0 & 0
		\end{pmatrix}}
	\end{alignat}

Spiegelung
	\begin{equation}
		\underline{S} =
		\begin{pmatrix}
			1 & 0 & 0 \\
			0 & 1 & 0 \\
			0 & 0 & -1
		\end{pmatrix}
	\end{equation}


\subsection{Bogenlänge und Zerlegung der Beschleunigung}
Die Bogenlänge ist der Betrag des infinitesimalen Wegelements
	\begin{equation}
		\diff s := \sqrt{\diff\vec{x} \cdot \diff\vec{x}} = |\vec{v}| \cdot \diff t
	\end{equation}

Tangenten-, Normalen- und Binomialvektor sind
	\begin{align*}
		\vec{e}_T &= \frac{\diff\vec{x}}{|\diff\vec{x}|} \equiv \frac{\diff\vec{x}}{\diff s} \\
		\vec{e}_N &= \frac{\diff \vec{e}_T}{\diff s} \frac{1}{|\frac{\diff \vec{e}_T}{\diff s}|} \\
		\vec{e}_B &= \vec{e}_T \times \vec{e}_N
	\end{align*}

Betrachtet man die Bahnkurve auf infinitesimaler Ebene bereits als gekrümmte Bahn, lässt sich die Beschleunigung in Tangential- und Normalbeschleunigung wie folgt zerlegen
	\begin{equation}
		\ddot{\vec{a}} = \dot{x} \cdot \vec{e}_T + \frac{\dot{x}^2}{R} \cdot \vec{e}_N
	\end{equation}
	
mit dem Radius des Schmiegekreises $R$.


\newpage


\subsection{Wegintegrale}
Um Wegintegrale zu berechnen, geht man wie folgt vor
\begin{enumerate}
	\item Parametrisierung der Raumkurve
		$$\vec{x} = \vec{x}(\tau),\qquad \tau\in [\tau_0,\tau_1]$$
	s.d. $\vec{x}(\tau_0) = P_0$ und $\vec{x}(\tau_1) = P_1$.
	
	\item Formuliere das Differential
		$$\diff\vec{x} = \frac{\diff \vec{x}}{\diff \tau} \diff\tau = \lr{ \frac{\diff x_1}{\diff \tau}\vec{e}_1 + \frac{\diff x_2}{\diff \tau}\vec{e}_2 + \frac{\diff x_3}{\diff \tau}\vec{e}_3 } \diff\tau$$
		
	\item Formuliere die Kraft
		$$\vec{K}(\vec{x}) = \vec{K}(\vec{x}(\tau))$$
	
	\item Einsetzen in das Integral
		$$A = \int_c \vec{K} \diff\vec{x} = \int_{\tau_0}^{\tau_1} \diff\tau \lr{ K_1(\tau)\cdot \frac{\diff x_1}{\diff \tau}\vec{e}_1 + K_2(\tau)\cdot \frac{\diff x_2}{\diff \tau}\vec{e}_2 + K_3(\tau)\cdot \frac{\diff x_3}{\diff \tau}\vec{e}_3 }$$
\end{enumerate}

\vspace{0.5cm}

{\color{blue}
\hrule
	\fbox{Beispiel}\\
	ges: Wegintegral im Kraftfeld $\vec{K} = (axy^3, bx^2y^2, cz^3)^T$ von $P_0 = (1,0,0)$ nach $P_1 = (0,1,1)$ auf einem schraubenförmigen Weg\\
	\\
	1) Parametrisierung der Raumkurve
		$$\vec{x}(\tau) = 
		\begin{pmatrix}
			\cos\tau \\ \sin\tau \\ \tau \frac{2}{\pi}
		\end{pmatrix},\qquad \tau\in [0,\frac{\pi}{2}]$$
	
	2) Formuliere das Differential
		$$\diff\vec{x} =
		\begin{pmatrix}
			-\sin\tau \\ \cos\tau \\ \frac{2}{\pi}
		\end{pmatrix} \diff\tau$$
	
	3) Formuliere die Kraft
		$$\vec{K} = 
		\begin{pmatrix}
			a\cos\tau\sin^3\tau \\
			b\cos^2\tau\sin^2\tau \\
			c \lr{\frac{2\tau}{\pi}}^3
		\end{pmatrix}$$
	
	4) Einsetzen in das Integral
		\begin{align*}
			A &= \int_0^{\pi/2} \diff\tau \begin{pmatrix}
			a\cos\tau\sin^3\tau \\
			b\cos^2\tau\sin^2\tau \\
			c \lr{\frac{2\tau}{\pi}}^3
		\end{pmatrix} \begin{pmatrix}
			-\sin\tau \\ \cos\tau \\ \frac{2}{\pi}
		\end{pmatrix} = \\
			 &= \int_0^{\pi/2} \diff\tau \lreck{ -a \frac{\diff}{\diff \tau} \frac{\sin^5\tau}{5}  + b \frac{\diff}{\diff\tau} \lr{ \frac{\sin^3\tau}{3} - \frac{\sin^5\tau}{5} } + c \lr{\frac{2}{\pi}}^4 \tau^3 } = -\frac{a}{5} + \frac{2b}{15} + \frac{c}{4}
		\end{align*}
\hrule
}


\newpage


\subsection{Legendre-Transformation}
Die Legendre-Transformation für eine Funktion, die von einer Variablen abhängt ist
	\begin{equation}
		g(p) = f(x) - \frac{\diff f}{\diff x}\cdot x
	\end{equation}

Für mehrdimensionale Funktionen gibt es die Möglichkeit eine oder mehrere Variablen zu tauschen.
	\begin{align*}
		g(p_1, x_2, ..., x_n) &= f(x_1,...,x_n) - \lr{ \frac{\partial f}{\partial x_1} }_{x_2,...,x_n} \cdot x_1 \\
		g(p_1, ..., p_n) &= f(x_1,...,x_n) - \sum_{i} \lr{ \frac{\partial f}{\partial x_i} }_{x_j,\ j\neq i} \cdot x_i
	\end{align*}

\textbf{Algorithmus}
\begin{enumerate}
	\item Notiere die alte Funktion $f$ und die neue Funktion $g$ auf
	\item Notiere die zu tauschende Variable $x_i$
	\item Berechne und notiere die neue Variable $p_i = \lr{ \frac{\partial f}{\partial x_i} }_{x_j}$
	\item Setze in die entsprechende Formel ein
\end{enumerate}

\vspace{0.5cm}

{\color{blue}
\hrule
	\fbox{Beispiel}\\
	\textcolor{red}{ERGÄNZEN}
\hrule
}


\newpage


\section{Formalismen der Analytischen Mechanik}
\subsection{Lagrange-Gleichungen 1. Art}
	\begin{align}
		m\ddot{\vec{x}} &= \vec{K} + \vec{Z} \equiv \vec{K} + \lambda \cdot \vec{\nabla}\varphi \\
		\varphi(\vec{x},t) &= \text{konst}.
	\end{align}

\textbf{Lösungsstrategie für kartesische Koordinaten}
	\begin{enumerate}
		\item Umstellen nach einer Koordinate
			\begin{align*}
				\varphi(x,y,z,t) &= \text{konst} \\
				\Leftrightarrow z &= f(x,y,t)
			\end{align*}
			
		\item Lagrange-Multiplikator mit eliminierter Koordinate\\
		Betrachte die Bewegungsgl. der zu eliminierenden Koordinate (hier $z$), setze $z$ und dessen Ableitungen ein und stelle nach $\lambda$ um.
			\begin{align*}
				m\ddot{z} &= K_z + \lambda(x,y,z,t) \\
				\Leftrightarrow \lambda(x,y,\dot{x},\dot{y},\ddot{x},\ddot{y},t) &= ...			
			\end{align*}
			
		\item Setze $z,\lambda$ in die Bewegungsgln. für $x,y$ ein\\
		Dann verbleiben
			\begin{align*}
				m\ddot{x} &= K_x(x,y,\dot{x},\dot{y},t) + \lambda(x,y,\dot{x},\dot{y},\ddot{x},\ddot{y},t) \cdot \frac{\partial\varphi}{\partial x}(x,y,z,t)\\
				m\ddot{y} &= K_y(x,y,\dot{x},\dot{y},t) + \lambda(x,y,\dot{x},\dot{y},\ddot{x},\ddot{y},t) \cdot \frac{\partial\varphi}{\partial y}(x,y,z,t)		
			\end{align*}
		
		\item Integration liefert $x(t),y(t)$ und somit $z(t),\lambda(t)$
	\end{enumerate}


\newpage


{\color{blue}
\hrule
	\fbox{Beispiel}\\
	geg: Schiefe Ebene mit konstantem Neigungswinkel $\alpha$\\
	ges: $\vec{x}(t)$\\
	\\
	Es gilt die Zwangsbedingung
		$$\varphi(x,y,z,t) = \frac{z}{x} = \tan\alpha = \text{konst.}$$
	
	Somit sind die Lagrange'schen Bewegungsgleichungen 1. Art:
		\begin{align*}
			m\ddot{x} &= -\lambda\tan\alpha \\
			m\ddot{y} &= 0 \\
			m\ddot{z} &= -mg + \lambda
		\end{align*}
	
	Verwende nun die im Kapitel gegebene Lösungsstrategie zur Lösung des Gleichungssystems:\\
	1) Umstellen von $\varphi = \text{konst.}$ nach $z$ und bestimme deren Ableitungen
		\begin{align*}
			\varphi &= z - x\tan\varphi = 0 \\
			\Leftrightarrow z &= x\tan\alpha \\
			\Rightarrow \dot{z} &= \dot{x}\tan\alpha \\
			\Rightarrow \ddot{z} &= \ddot{x}\tan\alpha
		\end{align*}
	
	2) Einsetzen in die Bew.gl. der zu eliminierenden Koordinate
		\begin{align*}
			m\ddot{z} &= -mg + \lambda \\
			\Rightarrow m (\ddot{x}\tan\alpha) &= -mg + \lambda \\
			\Leftrightarrow \lambda &= m\ddot{x}\tan\alpha + mg
		\end{align*}
	
	3) Einsetzen von $\lambda$ in die Bew.gl. von $x$
		\begin{align*}
			\cancel{m}\ddot{x} &= -(\cancel{m}\ddot{x}\tan\alpha + \cancel{m}g) \tan\alpha \\
			\Leftrightarrow \ddot{x} \underbrace{\lr{ \tan\alpha + \frac{1}{\tan\alpha}} }_{= \frac{\sin\alpha}{\cos\alpha} + \frac{\cos\alpha}{\sin\alpha} = \frac{1}{\cos\alpha\sin\alpha}} &= -g \\
			\Leftrightarrow \ddot{x} = -g\cos\alpha\sin\alpha
		\end{align*}
		
	4) Integration von $\ddot{x}$, um $x(t)$ und somit $z(t)$ zu erhalten
		\begin{align*}
			\stackrel{\comment{\int}}{\Rightarrow} x(t) &= -\frac{1}{2} g \cos\alpha\sin\alpha + v_0\cdot t + x_0 \\
			\Rightarrow z(t) &= x\cdot \tan\alpha \stackrel{\comment{x(t)}}{=} - \frac{1}{2} g \sin^2\alpha\cdot t^2 + (v_0\cdot t + x_0) \tan\alpha
		\end{align*}


\newpage


	\begin{figure}[h]
		\centering
		\includegraphics[width=0.6\textwidth]{lagrange 2 art schiefe ebene.pdf}
		\caption{Bewegung eines Massenpunkts entlang einer schiefen Ebene unter dem Winkel $\alpha=20^\circ$, Anfangsgeschwindigkeit $v_0 = 15\,\,\text{m/s}$ und Anfangsposition $x_0 = 15\,\text{m}$}
	\end{figure}

\hrule
}


\newpage


{\color{violet}
\hrule
	\fbox{Herleitung}\\
	\textcolor{red}{ERGÄNZEN}


\hrule
}


\newpage


\subsection{Lagrange-Gleichungen 2. Art}
\subsubsection{Einteilchensystem}
Bei $r$ holonomen Zwangsbedingungen
	$$\varphi_\nu = \text{konst},\qquad \nu = 1,..,r$$

besitzt der Massenpunkt
	$$f = 3-r$$

Freiheitsgrade. Dann werden nur $f$ geeignet gewählte Koordinaten $q_1,...,q_f$ benötigt, um die Bewegung des Körpers zu beschreiben. \textcolor{gruen}{Die restlichen Koordinaten werden mit den Zwangsbedingungen gleichgesetzt $q_{f+1}=\varphi_1,q_{f+2}=\varphi_2,...$}\\
\\
Stelle dann die Lagrange-Funktion $L = T - V$ auf (die nur von den geeignet gewählten Koordinaten $q_k$ abhängt)
	\begin{align}
		T &= \frac{m}{2}(\dot{x}_1^2 + \dot{x}_2^2 + \dot{x}_3^2) \\
		V &= -\int \vec{F}\,\diff \vec{x}
	\end{align}
		
und bestimme die Bewegungsgln. mittels
	$$\frac{\diff}{\diff t} \frac{\partial L}{\partial\dot{q}_k} - \frac{\partial L}{\partial q_k} = 0,\qquad k=1,...,f$$


\newpage


{\color{blue}
\hrule
	\fbox{Beispiel}\\
	geg: Schiefe Ebene mit konstantem Neigungswinkel $\alpha$\\
	\\
	Es gilt die Zwangsbedingung 
		$$\varphi(x,y,z,t) = \frac{z}{x} = \tan\alpha = \text{konst.}$$

	Somit ist der Freiheitsgrad $f = 3-1 = 2$. Es genügen also zwei Koordinaten, um das Problem zu beschreiben. Wähle die Polarkoordinaten, d.h. $q_1 = d, q_2 = y$ und
		\begin{align*}
			x &= d\cdot \cos\alpha \\
			z &= d\cdot \sin\alpha
		\end{align*}

	Bestimme die Lagrange-Funktion.\\
	Stelle dafür die kinetische Energie und das Potential auf, die nur von den geeignet gewählten Koordinaten $q_k$ abhängen dürfen.
		\begin{align*}
			T &= \frac{m}{2}(\dot{x}^2 + \dot{y}^2 + \dot{z}^2) = \frac{m}{2} \lr{ \dot{d}^2\cos^2\alpha + \dot{y}^2 + \dot{d}^2\sin\alpha} = \frac{m}{2} \lr{ \dot{d}^2 + \dot{y}^2 } \\
			F &= -mg \\
			\Rightarrow V &= mgz \stackrel{\comment{z}}{=} mg\cdot d\sin\alpha \\
			\Rightarrow L &= T - V = \frac{m}{2} \lr{ \dot{d}^2 + \dot{y}^2 } - mg\cdot d\sin\alpha
		\end{align*}

	Die Bewegungsgln. sind somit
		\begin{align*}
			\frac{\diff}{\diff t} \frac{\partial L}{\partial\dot{q}_k} - \frac{\partial L}{\partial q_k} &= 0 \\
		\intertext{\fbox{$d$}}
			\frac{\diff}{\diff t} m\dot{d} + mg\sin\alpha &= 0 \\
			\Leftrightarrow \ddot{d} &= -mg\sin\alpha \\
			\stackrel{\comment{\int}}{\Rightarrow} d(t) &= -\frac{1}{2}g\sin\alpha t^2 + \dot{d}_0 t + d_0
		\intertext{\fbox{$y$}}
			\frac{\diff}{\diff t}m\ddot{y} &= 0 \\
			\stackrel{\comment{\int}}{\Rightarrow} y(t) &= \dot{y}_0 t + y_0
		\end{align*}
	
\hrule
}


\newpage


{\color{violet}
\hrule
	\fbox{Herleitung}\\
	\textcolor{red}{ERGÄNZEN}




\hrule
}


\newpage


\subsubsection{Mehrteilchensystem}
Für Systeme mit mehreren Teilchen gilt dieselbe Lagrange-Gleichung, jedoch werden kinetische Energie und Potential wie folgt formuliert:
	\begin{align}
		T &= \sum_{j=1}^N \sum_{i=1}^3 \frac{m_i}{2} \lr{ \dot{x}_i^{(j)} }^2 \\
		 V &= \sum_{j=1}^N V^{(j)},
	\end{align}

wobei der hochgestellte Index $(j)$ das $j$-te Teilchen kennzeichnet.


\newpage


{\color{blue}
\hrule
	\fbox{Beispiel}\\
	geg: Doppelpendel in der x-y-Ebene, die Gravitationskraft zeige in y-Richtung\\
	\\
	Es gelten die Zwangsbedingungen
		\begin{align*}
			z^{(1)} &= 0 \\
			l_1 &= |\vec{x}^{(1)}| = \sqrt{ (x^{(1)})^2 + y^{(1)})^2 } \\
			z^{(2)} &= 0 \\
			l_2 &= |\vec{x}^{(2)} - \vec{x}^{(1)}|
		\end{align*}
	
	Somit gibt es
		$$f = 3N - 4 = 3\cdot 2 - 4 = 2$$

	Freiheitsgrade.\\
	Wähle die Winkel $q_1 = \varphi_1, q_2 = \varphi_2$ als Koordinaten, die die Bewegung beschreiben.\\
	Bestimme die Lagrange-Funktion.
		\begin{align*}
			T &= \frac{m_1}{2} \lreck{ (\dot{x}_1^{(1)})^2 + (\dot{x}_2^{(1)})^2  } + \frac{m_2}{2} \lreck{ (\dot{x}_1^{(2)})^2 + (\dot{x}_2^{(2)})^2 } = \\
			 &= \frac{m_1}{2}\lreck{ \lr{ \frac{\diff}{\diff t} l_1\sin\varphi_1 }^2 + \lr{ \frac{\diff}{\diff t} l_1\cos\varphi_2 }^2 } + \frac{m_2}{2} \left\{ \lreck{ \frac{\diff}{\diff t} (l_1\sin\varphi_1 + l_2\sin\varphi_2) }^2 + \right.\\
			  &+ \left.\lreck{ \frac{\diff}{\diff t} (l_1\cos\varphi_1 + l_2\cos\varphi_2) }^2 \right\} = \\
			  &=...= \frac{m_1}{2} l_1^2\dot{\varphi}^2 + \frac{m_2}{2} (l_1^2\dot{\varphi}_1^2 + l_2^2\dot{\varphi}_2^2 + 2l_1l_2\dot{\varphi}_1\dot{\varphi}_2 \cos(\varphi_1 - \varphi_2)) \\
			  V &= V^{(1)} + V^{(2)} = -m_1 g y^{(1)} - m_2 g y^{(2)} = \\
			  &= -m_1 g l_1\cos\varphi_1 - m_2 g (l_1\cos\varphi_1 + l_2\cos\varphi_2) = \\
			  &= -(m_1 + m_2) g l_1\cos\varphi_1 - m_2 g l_2\cos\varphi_2 \\
			  L &= T-V = \lr{ \frac{m_1 + m_2}{2} l_1^2\dot{\varphi}_1^2 + \frac{m_2}{2} l_2^2\dot{\varphi}_2^2 + m_2l_1l_2\dot{\varphi}_1\dot{\varphi}_2 \cos(\varphi_1 - \varphi_2)) } + \\
			   &+ \lr{ (m_1 + m_2) g l_1\cos\varphi_1 + m_2 g l_2\cos\varphi_2 }
		\end{align*}

	\fbox{$\varphi_1$}
		\begin{align*}
			\frac{\diff}{\diff t} \frac{\partial L}{\partial\dot{\varphi}_1} - \frac{\partial L}{\partial\varphi_1} &= 0 \\
			\Leftrightarrow \frac{\diff}{\diff t} \lr{(m_1 + m_2)l_1^2\dot{\varphi}_1 + 2m_2 l_1 l_2 \dot{\varphi}_2\cos(\varphi_1 - \varphi_2) } - \lr{ -m_2l_1l_2\dot{\varphi}_1^2\dot{\varphi}_2^2 \sin(\varphi_1 - \varphi_2)} &= 0 \\
			\Leftrightarrow ...
		\end{align*}

	\fbox{$\varphi_2$}
		\begin{align*}
			\frac{\diff}{\diff t} \frac{\partial L}{\partial\dot{\varphi}_2} - \frac{\partial L}{\partial\varphi_2} &= 0 \\
			\Leftrightarrow ...
		\end{align*}
	
	Numerisches Lösen ergibt
	\begin{figure}[h]
		\centering
		\includegraphics[width=0.6\textwidth]{Grafiken/doppelpendel.pdf}
	\end{figure}
		
\hrule
}


\newpage


\subsection{Hamilton-Formalismus}
	\begin{enumerate}
		\item Stelle die Lagrange-Funktion auf:
			$$L(q_1, ..., q_f, \dot{q}_1, ..., \dot{q}_f, t)$$
		
		\item Bilde die kanonisch konjugierten Impulse zu den $q_a$:
			$$p_a = \frac{\partial L}{\partial \dot{q}_a}$$

		\item Stelle die Hamilton-Funktion auf:
			\begin{equation}
				H(q_1, ..., q_f, p_1, ..., p_f, t) = \sum_{b = 1}^f p_b \cdot \dot{q}_b - L(q_1, ..., q_f, \dot{q}_1, ..., \dot{q}_f, t)
			\end{equation}
		
		\item Berechne die Hamilton'schen Bewegungsgleichungen
			\begin{align}
				\frac{\partial H}{\partial p_a} &= \dot{q}_a \\
				\frac{\partial H}{\partial q_a} &= -\dot{p}_a,
			\end{align}
		
		mit $a = 1, ..., f$.
	\end{enumerate}

\vspace{0.5cm}

{\color{blue}
\hrule
	\fbox{Beispiel}\\





\hrule
}


\newpage


{\color{violet}
\hrule
	\fbox{Herleitung}\\
	\textcolor{red}{ERGÄNZEN}

\hrule
}


\newpage


\subsection{Routh'sche Funktion}


\newpage


\subsection{Hamilton-Jacobi-Theorie}


\newpage


\section{Prinzip der kleinsten Wirkung}


\newpage


\section{Noether-Theorem}


\newpage


\section{Zweikörperproblem}
\subsection{Finite und infinite Bewegung am Diagramm}
Aus der Energieerhaltung und der Tatsache, dass die Bewegung in der Ebene statt findet ($x$-$y$-Ebene) folgt
	\begin{equation}
		E - \frac{l^2}{2mr^2} - V(r) \equiv E - V_\text{eff}(r) \geq 0
	\end{equation}

\begin{table}[h]
	\centering
	\begin{tabular}{l|l|l}
		\hline\hline
			$E < E_0$ & finite Bewegung & $r\in [ r_\text{min}, r_\text{max} ]$ \\
		\hline
			$E > E_0$ & infinite Bewegung & $r\in [ r_\text{min}, \infty)$ \\
		\hline\hline
	\end{tabular}
\end{table}


\subsection{Bahnkurve durch Integration}
Die Bahnkurve $r(\varphi)$ oder $\varphi(r)$ lässt sich mittels Integration bestimmen
	\begin{equation}
		\varphi - \varphi_0 = l \int_{r_0}^r \frac{1}{r^2} \frac{\diff r'}{\sqrt{ 2m(E - V(r')) - \frac{l^2}{r^2} }}
	\end{equation}







\subsection{Wirkungsquerschnitt}


\newpage


\section{Starrer Körper}
Für einen starren Körper gilt die Lagrange-Funktion
	\begin{align}
		L &= T - V = T_\text{Trans} + T_\text{Rot} - V = \frac{M}{2}\vec{V}^2 + \frac{1}{2} \vec{\omega}^T \underline{T} \vec{\omega} - V(X, Y, Z, \varphi, \theta, \psi)
	\end{align}


\subsection{Trägheitstensor $\underline{T}$}
Der Trägheitstensor taucht im Rotationsanteil der kinetischen Energie eines starren Körpers auf und lautet für den diskreten bzw. kontinuierlichen Fall
	\begin{align}
		\underline{T} &=
		\begin{pmatrix}
			\sum m (y^2 + z^2) & - \sum m x y & -\sum m x z \\
			- \sum m x y & \sum m (x^2 + z^2) & -\sum m y z \\
			-\sum m z x & -\sum m y z & \sum m (x^2 + y^2)
		\end{pmatrix} = \\
		&= 
		\begin{pmatrix}
			\int \diff V\, \rho (y^2 + z^2) & - \int \diff V\, \rho x y & -\int \diff V\, \rho x z \\
			- \int \diff V\, \rho x y & \int \diff V\, \rho (x^2 + z^2) & -\int \diff V\, \rho y z \\
			-\int \diff V\, \rho z x & -\int \diff V\, \rho y z & \int \diff V\, \rho (x^2 + y^2)
		\end{pmatrix}
	\end{align}

Für ein ausgewähltes körperfestes Koordinatensystem werden die Nebendiagonalelemente null und der Trägheitstensor somit diagonal
	\begin{align}
		\underline{T} &=
		\begin{pmatrix}
			I_1 & 0 & 0 \\
			0 & I_2 & 0 \\
			0 & 0 & I_3
		\end{pmatrix} =
		\begin{pmatrix}
			\sum m (y^2 + z^2) & 0 & 0 \\
			0 & \sum m (x^2 + z^2) & 0 \\
			0 & 0 & \sum m (x^2 + y^2)
		\end{pmatrix} = \\
		&= 
		\begin{pmatrix}
			\int \diff V\, \rho (y^2 + z^2) & 0 & 0 \\
			0 & \int \diff V\, \rho (x^2 + z^2) & 0 \\
			0 & 0 & \int \diff V\, \rho (x^2 + y^2)
		\end{pmatrix}
	\end{align}

\fbox{
\parbox{\textwidth}{
	Für hochsymmetrische Körper (z. B. Würfel, Kugel, Zylinder) liegen die Hauptachsen oft entlang der Symmetrieachsen.
}
}

Dann ist der rotationsanteil der kinetischen Energie
	\begin{equation}
		T_\text{Rot} = \frac{1}{2} \lr{ I_1\omega_1^2 + I_2\omega_2^2 + I_3\omega_3^2 }
	\end{equation}


\vspace{0.5cm}


{\color{blue}
\hrule
	\fbox{Beispiel}\\
	ges: Trägheitstensor eines homogenen Quaders\\
	\\
	Verwende das Koordinatensystem, das entlang der Symmetrieachsen des Quaders liegt, d.h. $\vec{e}_{x,y,z}$. Dann sind die Hauptträgheitsmomente einfach $I_{1,2,3}$.
		\begin{align*}
			I_1 &= \int \diff V\, \rho (y^2 + z^2) \stackrel{\comment{\rho = \frac{M}{V}}}{=} \frac{M}{V} \int_{-a/2}^{a/2} \diff x \int_{-b/2}^{b/2} \diff y \int_{-c/2}^{c/2} \diff z\, (y^2 + z^2) = \\
			&= \frac{M}{abc} \lreck{ ac\cdot \frac{1}{3} \lr{ (b/2)^3 - (-b/2)^3 } + ab\cdot \frac{1}{3} \lr{ (c/2)^3 - (-c/2)^3 } } = \\
			&= \frac{M}{abc} \frac{2}{3} \frac{1}{2^3} \lr{ ac b^3 + ab c^3} = \frac{M}{12} (b^2 + c^2)
		\end{align*}
\hrule
}


\subsection{Intrinsischer Drehimpuls $\vec{M}$}


\subsection{Drehmoment $\vec{D}$}


\subsection{Euler'sche Gleichungen}


\newpage


\section*{Appendix A: Rotierendes Bezugssystem um einen festen Punkt}
Sei $\Sigma$ das ruhende System. Ein Teilchen, das sich im rotierenden System $\Sigma'$ bewegt, folgt der Bewegungsgl.
	\begin{equation}
		m \frac{\diff'^2\vec{x}'}{\diff t} = \underbrace{\frac{\diff^2\vec{x}}{\diff t^2}}_{\color{gruen}=\vec{F}_{\text{ext}}} - 2m\vec{\omega} \times \frac{\diff'\vec{x}}{\diff t} - m\vec{\omega}\times(\vec{\omega}\times\vec{x}') - m\dot{\vec{\omega}}\times\vec{x}'
	\end{equation}

mit der neu definierten Ableitung $\frac{\diff'\vec{x}}{\diff t} = \sum_{i=1}^3 \dot{x}_i(t) \cdot \vec{e}_i(t)$ \textcolor{gruen}{(d.h. die Geschwindigkeit/Beschleunigung aus Sicht des rotierenden Systems)}, der Drehachse $\vec{\omega}$ und der Kreisfrequenz, mit der Sich das System dreht $\omega$.\\
\\
{\color{blue}
\hrule
	\fbox{Beispiel}\\




\hrule
}


\newpage


\section*{Appendix B: Umrechnung der Newton'schen Gleichungen}
Betrachte die Koordinatentransformation $\{x_1,x_2,x_3\} \to \{q_1,q_2,q_3\}$ durch die Funktionen
	$$x_i = f_i(q_1,q_2,q_3,t)$$

Die Newton'schen Bewegungsgleichungen im Satz der neuen Variablen $\{q_1,q_2,q_3\}$ ergeben sich über
	\begin{equation}
		Q_k \equiv \sum_{j=1}^3 \frac{\partial f_j}{\partial q_k} K_j = \frac{\diff}{\diff t} \frac{\partial}{\partial \dot{q}_k} T(q,\dot{q},t) - \frac{\partial}{\partial q_k} T(q,\dot{q},t)
	\end{equation}

mit der verallgemeinterten kinetischen Energie 
	$$T(q,\dot{q},t) = \frac{m}{2}\sum_{j=1}^3 \lr{ \sum_{i=1}^3 \frac{\partial f_j(q,t)}{\partial q_i} \dot{q}_i + \frac{\partial f_j(q,t)}{\partial t}  }^2 = \frac{m}{2} \sum_{i=1}^3 \dot{x}_i^2$$

\vspace{0.5cm}

{\color{blue}
\hrule
	\fbox{Beispiel}\\
	Newton'sche Bewegungsgln. in Zylinderkoordinaten
		\begin{alignat*}{2}
			x_1 &= f_1(q,t) =  q_1 \cos q_2 &&\equiv \rho \cos\varphi \\
			x_2 &= f_2(q,t) = q_1 \sin q_2 &&\equiv \rho \sin\varphi \\
			x_3 &= f_3(q,t) = q_3 &&\equiv z
		\end{alignat*}
	
	Die kinetische Energie ergibt sich aus
		\begin{align*}
			T &= \frac{m}{2}(\dot{x}_1^2 + \dot{x}_2^2 + \dot{x}_3^2) = \frac{m}{2} \lreck{ (\dot{\rho}\cos\varphi -\rho\dot{\varphi}\sin\varphi)^2 + (\dot{\rho}\sin\varphi + \rho\dot{\varphi}\cos\varphi)^2 + (\dot{z})^2 } = \\
			&= \frac{m}{2} (\dot{\rho}^2 + \rho^2\dot{\varphi}^2 + \dot{z}^2)
		\end{align*}

	Verwende nun die Gleichung für die generalisierten Kräfte\\
	1) $Q_1 \equiv Q_{\rho}$
		\begin{align*}
			& \frac{\partial}{\partial \rho} (\rho\cos\varphi) K_1 + \frac{\partial}{\partial \rho} (\rho\sin\varphi) K_2 + \frac{\partial}{\partial \rho} (z) K_3 =
			\begin{pmatrix}
				K_1\\K_2\\K_3
			\end{pmatrix}
			\begin{pmatrix}
				\cos\varphi\\ \sin\varphi\\ 0
			\end{pmatrix}
			= \vec{K} \vec{e}_r = K_r = \\
			=& \frac{\diff}{\diff t} \frac{\partial}{\partial \rho} \lr{ \frac{m}{2} (\dot{\rho}^2 + \rho^2\dot{\varphi}^2 + \dot{z}^2) } - \frac{\partial}{\partial \rho} \lr{ \frac{m}{2} (\dot{\rho}^2 + \rho^2\dot{\varphi}^2 + \dot{z}^2) } = m \ddot{\rho} - m\rho\dot{\varphi}^2 \\
			\Leftrightarrow & K_r =  m (\ddot{\rho} - \rho\dot{\varphi}^2)
		\end{align*}
	
	2) $Q_2 \equiv Q_{\varphi}$
		$$...\Rightarrow K_{\varphi} = m(\rho\ddot{\varphi} + 2\dot{\rho}\dot{\varphi})$$
		
	3) $Q_3 \equiv Q_z$
		$$...\Rightarrow K_z = m\ddot{z}$$
\hrule
}


\newpage


\section{Stichwortverzeichnis}
\begin{xltabular}{\linewidth}{p{6cm}|X}
	\hline\hline
		Begriff & Erklärung \\
	\hline
		Bogenlänge &  \\
	\hline
		Drehmoment &  \\
	\hline
		Euler'sche Gleichungen &  \\
	\hline
		Galilei-Transformation &  \\
	\hline
		generalisierte Kräfte &  \\
	\hline
		Hauptträgheitsachse &  \\
	\hline
		Hauptträgheitsmoment &  \\
	\hline
		homogene Funktion &  \\
	\hline
		Inertialsystem & Koordinatensystem, in dem sich ein Körper, auf den keine Kraft wirkt, in gleichförmiger Bewegung oder in Ruhe verbleibt \\
	\hline
		intrinsischer Drehimpuls &  \\
	\hline
		kanonische Transformationen &  Transformationen, die die Hamilton'schen Bewegungsgleichungen forminvariant lassen \\
	\hline
		kanonisch konjugierter Impuls &  \\
	\hline
		Kepler'sche Gesetze &  \\
	\hline
		Konfigurationsraum &  \\
	\hline
		konservative Kräfte & Kräfte, die aus einem Potential folgen
			$$\vec{K} = -\vec{\nabla}V$$ \\
	\hline
		Phasenraum &  \\
	\hline
		Prinzip der kleinsten Wirkung &  \\
	\hline
		Reparametrisierungsinvarianz &  \\
	\hline
		Schwerpunktsatz &  \\
	\hline
		Schwerpunktsystem &  \\
	\hline
		Trägheitstensor &  \\
	\hline
		verallgemeinerter Impuls &  \\
	\hline
		Virialsatz &  \\
	\hline
		Wegunabhängigkeit &  \\
	\hline
		Zwangsbedingung &  \\
	\hline
		Zwangskraft &  \\
	\hline
		zyklische Variable &  \\
	\hline\hline
\end{xltabular}


\end{document}




\section{Erhaltungsgrößen}
%\subsection{Energieerhaltung}
%Hängt das Potential $V$ nicht explizit von der Zeit ab, d.h. $\frac{\partial V}{\partial t} = 0$, ist die Energie
%	\begin{equation}
%		E := T + V = \text{konst.}
%	\end{equation}
%
%eine Erhaltungsgröße.
%
%\vspace{0.5cm}
%
%{\color{violet}
%\hrule
%	\fbox{Herleitung}\\
%	Sei $\frac{\partial V}{\partial t} = 0$ und $x_a(t)$ erfüllen die Newton'sche Bewegungsgl. \\
%	Untersuche die zeitliche Ableitung des Potentials.
%		\begin{align*}
%			\frac{\diff}{\diff t} V(x_1(t),...,x_{3N}(t)) &= \sum_{a=1}^{3N} \frac{\partial V}{\partial x_a} \frac{\partial x_a}{\partial t} + \underbrace{0}_{=\frac{\partial V}{\partial t}} = -\sum_{a=1}^{3N} F_a \dot{x}_a = - \sum_{a=1}^{3N} m \ddot{x}_a \dot{x}_a = -\frac{\diff}{\diff t} \sum_{a=1}^{3N} m \frac{1}{2} \dot{x}_a^2 = -\frac{\diff}{\diff t} T \\
%			\Leftrightarrow \frac{\diff}{\diff t} (T + V) &\equiv \frac{\diff}{\diff t} E = 0 \\
%			\Rightarrow E &= \text{konst.} \qquad\blacksquare
%		\end{align*}
%
%
%
%\hrule
%}
%
%
%
%\subsection{Impulserhaltung}
%Ist das Potential $V$ translationsinvariant, d.h. $V(\vec{x}^I,t) = V(\vec{x}^I + \vec{a},t)$, ist der Gesamtimpuls \textcolor{gruen}{(der die Summe über die Einzelimpulse aller Teilchen ist)}
%	\begin{equation}
%		\vec{P} := \sum_{I=1}^N \vec{p}^I = \sum_{I=1}^N m_I\dot{\vec{x}}^I
%	\end{equation}
%
%eine Erhaltungsgröße.
%
%\vspace{0.5cm}
%
%{\color{violet}
%\hrule
%	\fbox{Herleitung}\\
%	Sei $V(\vec{x}^I,t) \stackrel{!}{=} V(\vec{x}^I + \vec{a},t)$.
%		\begin{align*}
%			V(\vec{x}^I + \vec{a},t) &= T_{V,\vec{x}_0 = \vec{x}^I}(\vec{x}^I) = \\
%			&= V(\vec{x}^I,t) + \sum_{I=1}^N \left.\lr{ \frac{\partial V}{\partial x_1^{I}}((\cancel{x_1^I} + a_1) - \cancel{x_1^I}) + \frac{\partial V}{\partial x_2^{I}}((\cancel{x_2^I} + a_2) - \cancel{x_2^I}) + \frac{\partial V}{\partial x_3^{I}}((\cancel{x_3^I} + a_3) - \cancel{x_3^I}) }\right|_{\vec{x}^I} + \mathcal{O}(a^2) = \\
%			&= V(\vec{x}^I,t) + \sum_{i=1}^3 a_i \sum_{I=1}^N \frac{\partial V}{\partial x_i}(\vec{x}^I,t) + \mathcal{O}(a^2) \stackrel{!}{=} V(\vec{x}^I,t) \\
%		\Rightarrow 0 &\stackrel{!}{=} \sum_{i=1}^3 a_i \sum_{I=1}^N \frac{\partial V}{\partial x_i}(\vec{x}^I,t) \\
%		\Leftrightarrow 0 &\stackrel{!}{=} \sum_{I=1}^N \frac{\partial V}{\partial x_i}(\vec{x}^I,t) = -\sum_{I=1}^N K_i^I = -\sum_{I=1}^N m_I \ddot{\vec{x}}^I = \frac{\diff}{\diff t} \lr{ \sum_{I=1}^N m_I \dot{\vec{x}}^I } \equiv - \frac{\diff}{\diff t} \sum_{I=1}^N \vec{p}^I \equiv -\frac{\diff}{\diff t} \vec{P} \qquad\blacksquare
%		\end{align*}
%\hrule
%}
%
%
%\newpage
%
%
%\subsection{Drehimpulserhaltung}
%Ist das Potential $V$ rotationsinvariant, d.h. $V(\vec{x}^I,t) = V(\underline{A} \vec{x}^I,t)$, ist der Drehimpuls
%	\begin{equation}
%		\vec{L} = \sum_{I=1}^N m_I \vec{L}^I = \sum_{I=1}^N m_I^2 \vec{x}^I \times \dot{\vec{x}}^I
%	\end{equation}
%
%erhalten.
%
%\vspace{0.5cm}
%
%{\color{violet}
%\hrule
%	\fbox{Herleitung}\\
%	\textcolor{red}{ERGÄNZEN}
%
%
%
%\hrule
%}